% !TEX root = ../PhD Thesis.tex
\chapter{Side Project II - Charm: Comprehensive Hub for Alternative Regulatory Mapping of RBPs}
\label{chap:ChapterIII}

In Chapter 4, we focused on \glspl{rbp} known to be sequestered within \glspl{sg} and showed that functional alterations can be inferred from changes in their \gls{rna} binding profiles, even in the absence of differential expression. This establishes a framework to detect \gls{rbp} functional perturbations across diverse conditions using transcriptomic data alone. Compared to targeted approaches such as \gls{if} (for localisation changes) or western blotting (for post-translational modifications), this strategy enables broader and more scalable screening of candidate \glspl{rbp}.
\\
\\
Importantly, this approach does not distinguish the underlying mechanism driving functional changes, but it does provide an efficient first-pass method to identify \glspl{rbp} whose activity is likely perturbed, thereby prioritising targets for more in depth downstream analyses.
\\
\\
Our observations with \gls{tia1} serve as a proof of concept for this framework, demonstrating that \gls{sg}-associated sequestration can be linked to measurable functional consequences at the transcriptomic level. Building on this, we expanded the analysis beyond \gls{sg}-associated proteins to include all \glspl{rbp} with available ENCODE data\cite{the_encode_project_consortium_integrated_2012}, allowing for a more comprehensive and systematic assessment of \gls{rbp} functional modulation, not only on a \gls{sg} context, but in any physiological conditions in which their function is expected to altered.

\section{Materials and Methods}

\subsection{Differential Splicing Analysis}

We obtained \gls{rnaseq} data of gene silencing experiments from ENCODE\cite{the_encode_project_consortium_integrated_2012} targeting \glspl{rbp} for which corresponding \gls{eclip} data are also available from the same resource. The selected experiments were conducted in the HepG2 and K562 cell lines and employed either shRNA-mediated knockdown or \gls{crispr}-based perturbation.
\\
\\
For each cell line and target, we selected a single gene silencing experiment. Selection prioritised non-archived datasets; when multiple non-archived options were available, \gls{crispr}-based experiments were preferred. \gls{rnaseq} data for all replicates of the chosen experiment, along with their corresponding controls, were downloaded as FASTQ files.
\\
\\
Raw sequencing reads (FASTQ format) were first trimmed using Trimmomatic (v0.39)\cite{bolger_trimmomatic_2014} in \gls{pe} mode. Alternative splicing was quantified with VAST-TOOLS (v2.5.1)\cite{tapial_atlas_2017}, following the \texttt{align} (against human reference genome hg38) and \texttt{combine} steps.
\\
\\
Differential splicing analysis was them conducted using the betAS R package (v1.2.1)\cite{ascensao-ferreira_betas_2024}, with modifications to the default filtering procedure. By default, the \texttt{alternative\allowbreak Events} function excludes events in which at least one sample does not meet the \gls{psi} thresholds (1 $<=$ \gls{psi} $<=$ 99). However, applying these thresholds resulted in the inclusion of a large number of events that appeared constitutive, showing minimal variation rather than true alternative splicing.
\\
\\
To improve specificity for alternative splicing, we incorporated prior information from the VastDB MAIN PSI table\cite{tapial_atlas_2017}. For exon-skipping events, we computed the average \gls{psi} and variance across tissues and cell types. Events were classified as alternative if they satisfied at least one of the following empirical criteria: variance $>=$ 750, or average \gls{psi} between 12.5 and 87.5 (Figure 9.1). Only events meeting these criteria were retained for downstream filtering of the VAST-TOOLS output.

\begin{figure}[H]
	\centering
	\includegraphics[width=1\linewidth]{"images/Paper Images/eCLIPSE_FILT"}
	\caption[Alternative Events Selection]{\textbf{Alternative events selection.} Scatter plots of average PSI (X-Axis) and variance (Y-Axis) across tissues for exon skipping (Left) and intron retention (Right) events annotated in VastDB\cite{tapial_atlas_2017}. Both low (i.e., close to 0) and high (i.e., close to 1) average \gls{psi} values may indicate that an event is more cryptic (i.e., an event that is almost always constitutive, but, under specific conditions, such as failure of other splice sites, is alternative instead\cite{aldalaqan_cryptic_2022}) than alternatively regulated. Because these metrics are averaged across multiple tissues, low variance can reflect very tissue-specific regulation, such as an event consistently included in one tissue (e.g., neuronal) but excluded in others. For downstream analyses, only events with strong evidence of alternative inclusion were maintained. Dashed lines indicate the applied filtering thresholds: variance $>=$ 750 (horizontal) and PSI between 12.5 and 87.5 (vertical).}
	\label{fig:Figure6.1}
\end{figure}



\subsection{\textit{In Silico} Resources}

All analyses were performed on the workstation \textit{lobito} at the \gls{gimm} data center, which is equipped with 19 \gls{tb} of \gls{ssd} storage, 768 \gls{gb} of RAM, 72 \gls{cpu} threads (\gls{intel} Xeon
Gold 6254 @ 3.10 \gls{ghz}), and an integrated Matrox G200eW3 \gls{gpu}. \gls{bash} was used to run VAST-TOOLS\cite{tapial_atlas_2017}, FASTQC\cite{andrews_fastqc_2010}, and MultiQC\cite{ewels_multiqc_2016}. Downstream analyses were conducted in R (v4.1.2) using RStudio Server (v1.4.1717). The eCLIPSE tool includes steps that were conducted in Python 3 (v3.10.11), and in R (v4.1.2). CHARM was created as an open usage suite with RShiny (v1.11.1), and is planned to be hosted in the host lab's tool/app website.


\section{RBP Binding Analyses Pipeline (eCLIPSE)}

\gls{eclipse} is a computational pipeline developed to integrate \gls{eclip} data from the ENCODE project\cite{the_encode_project_consortium_integrated_2012} with splicing annotations from \gls{vastdb}\cite{tapial_atlas_2017}, generating \gls{rna} splicing maps. By aligning \gls{rbp} binding profiles to the genomic coordinates of annotated splicing events, \gls{eclipse} enables the systematic assessment of how individual \glspl{rbp} regulate alternative splicing, either directly (via their own binding) or indirectly (through downstream effects on other \glspl{rbp}). All components required to reproduce this pipeline are available on GitHub\footnote{at https://github.com/DiseaseTranscriptomicsLab/CHARM}, with the exception of large files that cannot be hosted publicly due to size constraints but are nonetheless available upon request.
\\
\\
A brief description of each relevant step is shared in the sections below.

\subsection{Splicing-Genome Creation}

A reference ``splicing-genome" was first created from VastDB annotations, for exon skipping, and for intron retention events.
\\
\linebreak
\noindent
Alternative exons shorter than 10 bp were excluded. For each exon skipping event, specific regions of interest (as exemplified on Figure 4.1)  were defined as follows:

\begin{verbatim}
	Upstream exon end: -50 bp to the exon boundary;
	Upstream intron start: intron boundary to +200 bp;
	Upstream intron end: -200 bp to the intron boundary;
	Alternative exon start: exon boundary to +50 bp;
	Alternative exon end: -50 bp to the exon boundary;
	Downstream intron start: intron boundary to +200 bp;
	Downstream intron end: -200 bp to the intron boundary;
	Downstream exon start: exon boundary to +50 bp.
\end{verbatim}
\noindent
Additional rules for handling short exons or introns ($<$100 bp or $<$400 bp, respectively) were defined.  When upstream or downstream introns were shorter than 400 bp, they were divided between the start and end regions, to ensure we maintained comparable regions of interest (i.e., near the splice sites) for these events. For even-length introns, the split was equal; for odd-length introns, the extra base was assigned to the end half. For alternative exons shorter than 100 bp, the same rule was applied. For alternative exons with alternative splice sites, the splice site that maximised exon length was chosen.
\\
\linebreak
\noindent
Intron retention events with introns shorter than 30 bp were excluded. For each intron retention event, specific regions of interest were defined as follows:

\begin{verbatim}
	Upstream exon end: -50 bp to the exon boundary;
	Intron start: intron boundary to +200 bp;
	Intron end: -200 bp to the intron boundary;
	Downstream exon start: exon boundary to +50 bp.
\end{verbatim}
\noindent 
Additional rules for handling short introns ($<$400 bp) were defined. These introns' length was divided between the start and end regions, to ensure we maintained comparable regions of interest for these events. For even-length introns, the split was equal; for odd-length introns, the extra base was assigned to the end region. 
\\
\\
This step was conducted in R (v4.1.2).

\subsection{eCLIP Samples}

One \gls{eclip} file per cell line for each available \gls{rbp} in \gls{encode}\cite{the_encode_project_consortium_integrated_2012} was considered. Processed \gls{hg38} \gls{bed} narrowPeak files from each \gls{rbp} labelled as ``default" by \gls{encode}\cite{the_encode_project_consortium_integrated_2012} were downloaded, and binding events with the values in the ``Signal" column greater than 3 were retained.

\subsection{Binding Metagenome Assembly}

\gls{rbp} binding peaks were intersected with the predefined splicing regions to extract the binding signal in \gls{hg38} coordinates. For each \gls{rbp}, each base in the original genome was coded as 0 (no binding) or 1 (binding present). To enable comparison across events of different lengths, each event was mapped to a standardised ``metagenome" coordinate system corresponding to the regions of interest (e.g., upstream exon, upstream intron, alternative exon, downstream intron, downstream exon, as defined in 9.2.1). Each splicing event was subsetted and rescaled to fixed metagenome coordinates (e.g., 1-50 corresponding to -50 nt to upstream exon end, 51-250 corresponding to upstream intron start to +200 nt, …), preserving the relative position of binding within the region. As stated above, because individual exons or introns can be shorter than the full reference \gls{roi} (i.e., exons $<$100 nt or introns $<$400 nt), some positions in the metagenome may not have corresponding \gls{hg38} coordinates with binding information. These positions are assigned \textit{\gls{na}} to indicate missing binding data. Binding information in metagenome coordinates were strand-specific. All conversions were implemented in Python3 (v3.10.11). Since Python indexing starts at 0, a +1 was added to Python defined coordinates, to ensure consistency with R defined coordinates, whose indexing starts at 1.

\subsection{RNA-Binding Maps}

betAS-derived\cite{ascensao-ferreira_betas_2024} differential splicing events were divided into categories as maintained, increased, or decreased, based on a chosen dPSI threshold.
\\
\\ 
For each of these categories, and per \gls{rbp}, normalised binding profiles were computed across the metagenome coordinates for the corresponding events. Specifically, for each nucleotide position in the metagenome, the binding signal was calculated to ensure a per-position signal that accounts for missing values due to shorter regions in individual events, which would otherwise be overrepresented:

\begin{equation}
	\text{Normalised binding} =
	\frac{\text{Sum of events with binding}}{\text{Total events} - \text{Sum of NAs}}
\end{equation}
\eqcaption{Normalised binding procedure}
\\
\noindent
For assessment of differential binding, a Chi-Squared Test was applied to each position using the following contingency matrix (Table 9.1): 


\begin{table}[H]
	\centering
	\caption[Contingency matrix used for Chi-Squared Test in eCLIPSE]{\textbf{Contingency matrix used for Chi-Squared Test in eCLIPSE}. For these tests, increased events, or decreased events were compared to maintained events alone.}
	\scriptsize
	
	\begin{tabular}{|p{6cm}|p{6cm}|}
		\hline
		Number of Increased/Decreased Events & Number of Maintained Events \\
		\hline
		Number of Increased/Decreased Events after Normalised Binding &
		Number of Maintained Events after Normalised Binding \\
		\hline
	\end{tabular}
	
	\label{tab:event_counts}
\end{table}

\noindent
The resulting p-values from the test are then log10-transformed and corrected for multiple testing using the \gls{bh} procedure\cite{benjamini_controlling_1995}, resulting in \gls{fdr}, which were multiplied by -1 for positions with decreased binding.
\\
\\
Finally, position-wise normalised binding profiles were plotted along the metagenome coordinates for each category, with the corresponding FDR values, or Chi-Squared Statistic values, shown in the bottom half of the plot. Taken together, these steps generate statistically assessed RBP binding maps for maintained, increased, and decreased exon-skipping events. This process was performed in R (v4.1.2).

\subsection{dPSI Threshold Optimisation}

For integration with CHARM's binding tab, additional optimisations were performed.
\\
\\
As described in Section 9.2.4, the categorisation of splicing events into increased, decreased, and maintained classes depends on the chosen \gls{dpsi} threshold, which defines the minimum magnitude of splicing change required to assign an event to the increased or decreased class. Because the optimal threshold may vary across \glspl{rbp} depending on binding strength and regulatory effect size, a systematic optimisation was performed prior to differential binding analysis.
\\
\\
For each \gls{rbp} \gls{kd}, thresholds ranging from 0.01 to 0.10 (in increments of 0.01) were evaluated across all available \gls{eclip} binding datasets (that is, the binding profiles of all other \glspl{rbp} as well as the \gls{kd} \gls{rbp} itself). At each threshold, splicing events were categorised as described in Section 9.2.4, and the \gls{rna} splicing map was computed along with position-wise \gls{fdr} values and Chi-Squared statistics. To obtain a single summary of differential binding strength across the metagenome, the signed \gls{fdr} and Chi-Squared values were summed in absolute terms across all 1000 metagenome positions, separately for the increased and decreased splicing classes. Here, signed refers to the convention described in Section 9.2.4, whereby values at positions with decreased binding relative to the maintained class are multiplied by $-1 $ prior to summation.
\\
\\
This procedure yielded four summary scores per threshold per \gls{kd}-\gls{eclip} pair: one for each combination of splicing category (increased or decreased) and metric (\gls{fdr} or Chi-Squared). The threshold maximising each of these four scores was retained as optimal for that combination. Consequently, up to four distinct optimal thresholds may be selected per \gls{kd}-\gls{eclip} pair. In addition, thresholds of 0.05 and 0.10 were always retained as standard reference values, irrespective of the optimisation outcome.

\section{From eCLIPSE to CHARM}

With the eCLIPSE pipeline established, we developed a systematic framework to analyse alternative splicing in the context of biological perturbations. Starting from gene silencing experiments, the pipeline enables quantification of splicing events as increased, decreased, or maintained, and integrates differential binding enrichment to infer functional alterations in \glspl{rbp}.
\\
\\
Using \gls{encode} \gls{shrna} datasets targeting \glspl{rbp}, we assembled a multi-layered database capturing gene expression, splicing, and protein–RNA binding. While integrative analyses across these layers have been previously reported\cite{rogalska_transcriptome-wide_2024}, our approach introduces a key extension: the ability to assess indirect binding effects, specifically how \gls{kd} of one \gls{rbp} putatively influences the binding profiles of others. This added dimension enables the exploration of cooperative and competitive interactions, facilitating the identification of candidate regulatory relationships and potential RBP complexes.
\\
\\
To support this analysis, we developed \gls{charm} as an open-access, modular suite implemented as an RShiny web application. Its design allows straightforward integration of additional datasets, including gene silencing experiments from other sources, \gls{eclip} profiles for additional RBPs, and alternative \gls{clip} technologies such as iCLIP2\cite{buchbender_improved_2020}. CHARM is intended as a hypothesis-generation platform, with results requiring downstream experimental validation.
\\
\\
This framework enables the identification of functional alterations in RBPs that are not evident at the expression level alone. Furthermore, it facilitates the detection of alternative splicing events, such as intron retention, that may generate neoantigens associated with clinical outcomes or serve as candidates for therapeutic targeting, an area of growing translational interest\cite{rogalska_regulation_2023}.
\\
\linebreak
\noindent
Thus, CHARM has two modes of usage (Figure 9.2): (1) Explore, in which users can investigate these layers of information derived from ENCODE's data\cite{the_encode_project_consortium_integrated_2012}; (2) Discovery, in which users can input their own expression or splicing data, to either compare with ENCODE's data\cite{the_encode_project_consortium_integrated_2012}, or to use the eCLIPSE pipeline to uncover differential binding.
\\

\begin{figure}[H]
	\centering
	\includegraphics[width=1\linewidth]{"images/Paper Images/CHARM_Home"}
	\caption[CHARM's Home Screen]{\textbf{CHARM's home screen.} Screenshot of the CHARM application home screen, illustrating the two available usage modes and their respective descriptions. Users can select a module by clicking the corresponding tab in the upper-left panel. The source code is accessible via the GitHub link located in the bottom-right corner.}
	\label{fig:Figure6.2}
\end{figure}
\noindent
CHARM was thus developed to integrate several layers (modules) of information pertaining to \gls{rbp}s, based on expression, splicing, and binding, reflecting how these are changed upon \gls{shrna} treatment, or other conditions of interest. An overview of CHARM's different modules is presented below.

\subsection{Expression}

Within the \textbf{Expression} tab, in \textbf{Explore Mode} (Figure 9.3A) users can explore the effects of \gls{rbp} \gls{kd} on gene expression. Analyses can be performed for HepG2 and K562 cell lines, either jointly or independently (Figure 9.3B). After selecting a \gls{rbp} (Figure 9.3C), the interface initially presents violin plots summarising the distributions of expression in control and \gls{kd} samples (Figure 9.3F), alongside a scatter plot of differential gene expression (X-Axis - Log2 Fold-Change; Y-Axis - B-Statistic) between \gls{kd} and control samples, allowing the user to compare the selected \gls{rbp} \gls{kd} with all other available \glspl{rbp} for reference (Figure 9.3G).
\\
\\
This comparative view is particularly relevant because some \gls{encode} \gls{kd} experiments show limited or inconsistent perturbation efficiency, with minimal or even increased expression of the targeted gene. In such cases, analyses remain available; however, a warning message is displayed (if the \gls{kd} has B-Statistic $<$ 0, or Log2 Fold-Change $>$ -0.5) to indicate potential data quality concerns.
\\
\\
Downstream results include a volcano plot of \gls{dgea} between \gls{kd} and control (Figure 9.3H) and an associated results table (Figure 9.3I). Both elements are interactive and implemented using Plotly (v4.11.0): selecting a gene in the table highlights the corresponding dot in the volcano plot, and vice-versa. The targeted \gls{rbp} is specifically highlighted in red.
\\
\\
\gls{gsea} results are presented below these plots (Figure 9.4A), showing enriched and depleted Hallmark gene sets\cite{liberzon_molecular_2015} along with a corresponding interactive table (Figure 9.4B).
\begin{figure}[H]
	\centering
	\includegraphics[width=1\linewidth]{"images/Paper Images/CHARM_Expression1"}
	\caption[CHARM's Expression Tab in Explore Mode, Part 1]{\textbf{CHARM's expression tab in Explore Mode, part 1.} (A) Selection of usage mode; in this example, Explore Mode is active. (B) Selection of data scope, allowing analysis of HepG2 and K562 cell lines either jointly or separately. This panel also provides access to the ``Similar RBPs" option. (C) RBP KD-centered analysis, where users select an RBP of interest to inspect what expression changes occur upon its \gls{kd} (D) Gene-centered analysis, where users select a gene of interest. Resulting visualisations display bar plots indicating the \glspl{rbp} that most strongly influence the expression of the selected gene (see Figure 9.5A). (E) Hallmark gene set–centered analysis, where users select a gene set. Resulting visualisations display bar plots showing the \glspl{rbp} most strongly associated with enrichment or depletion of the selected gene set (See Figure 9.5B). (F) Violin plots showing expression distributions of selected genes across relevant samples. When ``Both Cells" is selected, individual data points correspond to both HepG2 and K562 samples. Expression is reported as normalised Log2 CPM. (G) Scatter plot of differential gene expression (X-Axis - Log2 Fold-Change; Y-Axis - B-Statistic) between \gls{kd} and control samples, allowing the user to compare the selected \gls{rbp} \gls{kd} with all other available \glspl{rbp} for reference. (H) Volcano plot of differential gene expression (X-Axis, Log2 fold change; Y-Axis, B-statistic), between selected RBP's \gls{kd} and control samples. Genes are represented as circles; the targeted RBP is highlighted in red, while user-selected genes are highlighted in green. (I) Interactive DGEA results table corresponding to panel (H). Users can sort by metrics of interest, and selecting a row highlights the corresponding gene in the volcano plot.}
	\label{fig:Figure6.3}
\end{figure}

\begin{figure}[h]
	\centering
	\includegraphics[width=1\linewidth]{"images/Paper Images/CHARM_Expression1_2"}
	\caption[CHARM's Expression Tab in Explore Mode, Part 2]{\textbf{CHARM's Expression Tab in Explore Mode, Part 2.}  (A) Bar plot showing NES of GSEA t-statistics of differential gene expression between \gls{kd} and control samples for the selected RBP.  (B) Interactive GSEA results table corresponding to panel (A). Users can sort by metrics of interest.}
	\label{fig:Figure6.32}
\end{figure}
\noindent
The Expression tab also provides additional discovery-oriented submodes:

\begin{itemize}
\item Gene-centered submode (Figures 9.3D, 9.5A): Users can select a gene of interest and visualise, in a bar plot, the RBPs whose KD most strongly influence its expression in both positive and negative directions.
\item GSEA-centered submode (Figures 9.3E, 9.5B): Users can select a Hallmark gene set\cite{liberzon_molecular_2015} and evaluate which RBPs' KD are most strongly associated with its enrichment or depletion.
\end{itemize}


\begin{figure}[H]
	\centering
	\includegraphics[width=1\linewidth]{"images/Paper Images/CHARM_Expression2"}
	\caption[CHARM's Expression Tab in Explore Mode, Gene and Gene Set Centered submodes]{\textbf{CHARM's expression tab in Explore Mode, gene and gene set centered submodes} (A) Selection of gene for Gene-Centered mode, and corresponding bar plot, showcasing RBPs that lead to upregulation or downregulation of the selected gene, as given by Log2FC. (B) Selection of Hallmark Gene Set\cite{liberzon_molecular_2015} for Gene Set-Centered mode, and corresponding bar plot, showcasing RBPs that lead to stronger enrichment or depletion of the selected gene set, as given by NES.}
	\label{fig:Figure6.4}
\end{figure}
\noindent
Additionally, a ``Similar RBPs" option enables direct comparison between differential expression profiles associated with each RBP's KD (Figure 9.6). Visualisation is provided as scatter plots when a single comparison is selected, or bar plots for multiple \glspl{rbp}, like in Figure 9.7. Users can also identify the most and least similar \glspl{rbp} relative to a user's query (i.e., user selects RBP A, and is presented with RBPs whose KD is most similar, or dissimilar, to RBP A's KD.). Similarity can be defined based on gene expression correlations or GSEA Hallmark\cite{liberzon_molecular_2015} correlations (Figure 9.6C).
\begin{figure}[H]
	\centering
	\includegraphics[width=1\linewidth]{"images/Paper Images/CHARM_Expression3"}
	\caption[CHARM's Expression Tab in Explore Mode, Similar RBPs submode]{\textbf{CHARM's expression tab in explore mode, Similar RBPs option.} (A) Selection of the ``Similar RBPs" submode. (B) Selection of the reference RBP for similarity comparisons. (C) Selection of similarity metric; in this example, similarity is defined by gene expression correlation. (D) Selection of one or more RBPs for comparison with the reference RBP (B). If a single RBP is selected, results are displayed as a scatter plot (see panel F). If multiple RBPs are selected, results are shown as a bar plot of Spearman's correlation coefficient $\rho$, like in Figure 9.7. (E) Exploratory mode, where users can identify RBPs with the strongest positive or negative correlations relative to the reference RBP. Users may specify the number of RBPs to display, or select top positively and negatively correlated RBPs separately. Results are presented as bar plots ordered by decreasing absolute Spearman's $\rho$, like in Figure 9.7. (F) Scatter plot comparing differential expression t-statistics of the KDs of RBPs selected in (B) and (D). Spearman's $\rho$ and associated p-value are shown in the plot. Comparison given is from both cell lines analysed jointly, but user can scroll down to observe them analysed separately.}
	\label{fig:Figure6.5}
\end{figure}
\noindent
In Discovery Mode (Figure 9.7A), users may also upload their own differential gene expression dataset (Figure 9.7B). The input requires two columns: \gls{hgnc} gene symbols and a differential expression statistic (preferably t-statistic). Users can then compare their dataset against selected \glspl{rbp}'s KD using scatter or bar plot visualisations, or perform exploratory analyses to identify the most and least related RBPs. As in the ``Similar RBPs" mode, similarity between RBPs' KD can be assessed using either gene expression or GSEA-based correlations (Figure 9.7C).
\begin{figure}[H]
	\centering
	\includegraphics[width=1\linewidth]{"images/Paper Images/CHARM_Expression4"}
	\caption[CHARM's Expression Tab in Discovery Mode]{\textbf{CHARM's expression tab in Discovery Mode} (A) Selection of ``Discovery Mode". (B) File upload interface, with description of file format requirements above. (C) Selection of similarity metric; in this example, similarity is defined by gene expression correlation. (D) Selection of one or more RBPs for comparison with the user's uploaded file (B). If a single RBP is selected, results are displayed as a scatter plot, like in Figure 9.6. If multiple RBPs are selected, results are shown as a bar plot of Spearman's correlation coefficient $\rho$ (see panel F). (E) Exploratory Mode, where users can identify RBPs with the strongest positive or negative correlations relative to the user's uploaded file. Users may specify the number of RBPs to display, or select top positively and negatively correlated RBPs separately. Results are presented as bar plots ordered by decreasing absolute Spearman's $\rho$. (F) Bar plot of Spearman's correlations between the differential expression profile of each of the RBPs selected in (D) and that uploaded in (B). Comparison given is from both cell lines analysed jointly, but user can scroll down to observe them independently.}
	\label{fig:Figure6.6}
\end{figure}
\noindent

\subsection{Splicing}

Within the Splicing tab, in Explore Mode (Figure 9.8A), users can investigate how \gls{rbp} \gls{kd} reshapes splicing patterns, focusing specifically on exon skipping and intron retention events. Upon selection of an \gls{rbp} (Figure 9.8C), the interface summarises global splicing alterations (Figure 9.8E) and provides an interactive volcano plot displaying \gls{dpsi} (X-Axis) \gls{vs} the \gls{pdiff} (Y-axis) \cite{ascensao-ferreira_betas_2024}) (Figure 9.8G). Splicing events are identified by \gls{vastdb} IDs\cite{tapial_atlas_2017}, and results are accompanied by an interactive table for detailed inspection (Figure 9.8H).
\\
\\
Analyses can be performed on HepG2 and K562 cell lines either jointly or separately (Figure 9.8B). In contrast to the global view, an event-centered mode allows users to query a specific splicing event (Figure 9.8D) and visualise, via bar plots, which \glspl{rbp} KDs most strongly promote exon inclusion or exclusion (Figure 9.9).
\begin{figure}[H]
	\centering
	\includegraphics[width=1\linewidth]{"images/Paper Images/CHARM_Splicing1"}
	\caption[CHARM's Splicing Tab in Explore Mode, Part 1]{\textbf{CHARM's splicing tab in Explore Mode, part 1. } (A) Selection of usage mode; in this example, Explore Mode is active. (B) Selection of data scope, allowing analysis of HepG2 and K562 cell lines either jointly or separately. This panel also provides access to the ``Similar RBPs" option. (C) Selection of \gls{rbp} of interest. (D) Event-centered analysis, where users select an event of interest, as per VastDB IDs\cite{tapial_atlas_2017}. Resulting visualisations display bar plots indicating the \glspl{rbp} whose KD most strongly influence the alternative splicing of the selected event. (E) Violin plots showing global distribution of \gls{dpsi} of intron retention (IR) and exon skipping (EX) events. In the example shown, \textit{AQR} KD has a strong global effect in splicing efficiency, as its knockdown increases intron retention and alternative exon inclusion. (F) Scatter plot of differential gene expression (X-Axis - Log2 Fold-Change; Y-Axis - B-Statistic) between \gls{kd} and control samples, allowing the user to compare the selected \gls{rbp} \gls{kd} with all other available \glspl{rbp} for reference. (G) Volcano plot of differential splicing analysis (X-Axis, dPSI; Y-Axis, \gls{pdiff}\cite{ascensao-ferreira_betas_2024}) between selected RBP's KD and control samples. Events are represented as circles. (H) Interactive differential splicing results table corresponding to panel (G). Users can sort by metrics of interest, and selecting a row highlights the corresponding event in the volcano plot.}
	\label{fig:Figure6.7}
\end{figure}

\begin{figure}[H]
	\centering
	\includegraphics[width=1\linewidth]{"images/Paper Images/CHARM_Splicing1_2"}
	\caption[CHARM's Splicing Tab in Explore Mode, Part 2]{\textbf{CHARM's splicing tab in Explore Mode, part 2. } Bar plot of Event-Centered mode (as selected in Figure 9.8D), showcasing RBPs that lead to inclusion or exclusion of the selected alternative splicing event.}
	\label{fig:Figure6.72}
\end{figure}
\noindent
The “Similar RBPs” option enables comparative analysis of RBPs based on their splicing impact profiles (Figure 9.10A). When comparing a single pair of RBPs, results are shown as scatter plots of event-level agreement (Figure 9.10E); when multiple \glspl{rbp} are selected, similarity is summarised using bar plots based on correlation metrics, like in Figure 9.11. This option also supports exploratory queries to identify \glspl{rbp} with the most similar or dissimilar splicing signatures relative to a chosen reference (Figure 9.10D).
\begin{figure}[H]
	\centering
	\includegraphics[width=1\linewidth]{"images/Paper Images/CHARM_Splicing2"}
	\caption[CHARM's Splicing Tab in Explore Mode, Similar RBPs submode]{\textbf{CHARM's splicing tab in Explore Mode, Similar RBPs option.} (A) Selection of the ``Similar RBPs" submode. (B) Selection of the reference RBP for similarity comparisons. (C) Selection of one or more RBPs for comparison with the reference RBP (B). If a single RBP is selected, results are displayed as a scatter plot (see panel F). If multiple RBPs are selected, results are shown as a bar plot of Spearman's correlation coefficient $\rho$, like in Figure 9.11. (D) Exploratory Mode, where users can identify RBPs with the strongest positive or negative correlations relative to the reference RBP. Users may specify the number of RBPs to display, or select top positively and negatively correlated RBPs separately. Results are presented as bar plots ordered by decreasing absolute Spearman's $\rho$, like in Figure 9.11. (F) Scatter plot comparing differential splicing dPSIs of the KDs of RBPs selected in (B) and (C). Spearman's $\rho$ and associated p-value are shown in the plot. Comparison given is from both cell lines analysed jointly, but user can scroll down to observe them analysed separately.}
	\label{fig:Figure6.8}
\end{figure}
\noindent
In Discovery Mode (Figure 9.11A), users can upload an external differential splicing dataset (Figure 9.11B). To ensure compatibility, input tables must be preprocessed using the betAS application\cite{ascensao-ferreira_betas_2024}, which formats outputs derived from VAST-TOOLS\cite{tapial_atlas_2017} (and from other splicing quantification tools). Uploaded datasets can then be compared against reference \glspl{rbp} using the same visualisation framework as above (scatter or bar plots) (Figure 9.11C), or used to identify RBPs with concordant or opposing splicing effects (Figure 9.11D, F). Additionally, a plot of the global summary of exon skipping and intron retention alterations, is generated for the uploaded data (Figure 9.11E).
\begin{figure}[H]
	\centering
	\includegraphics[width=1\linewidth]{"images/Paper Images/CHARM_Splicing3"}
	\caption[CHARM's Splicing Tab in Discovery Mode]{\textbf{CHARM's splicing tab in Discovery Mode. } (A) Selection of ``Discovery Mode". (B) File upload interface, with description of file format requirements above. (C) Selection of one or more RBPs for comparison with the user's uploaded file (B). If a single RBP is selected, results are displayed as a scatter plot, like in Figure 9.6. If multiple RBPs are selected, results are shown as a bar plot of Spearman's correlation coefficient $\rho$ (see panel F). (D) Exploratory Mode, where users can identify RBPs with the strongest positive or negative correlations relative to the user's uploaded file. Users may specify the number of RBPs to display, or select top positively and negatively correlated RBPs separately. Results are presented as bar plots ordered by decreasing absolute Spearman's $\rho$. (E) Violin plots showing global distribution of \gls{dpsi} of intron retention (IR) and exon skipping (EX) events, in user's supplied data. In this case, uploaded data only contains EX events. (F) Bar plot of Spearman's correlations between the differential splicing profile of each of the RBPs selected in (C) and that uploaded in (B). Comparison given is from both cell lines analysed jointly, but user can scroll down to observe them independently.}
	\label{fig:Figure6.9}
\end{figure}
\noindent

\subsection{Binding}

Within the Binding tab, users can analyse putative changes in \gls{rbp}–\gls{rna} interaction patterns in the context of alternative splicing events, including exon skipping or intron retention. In this framework, \textit{differential binding} refers to differences in \gls{rbp} eCLIP signal distribution between groups of splicing events stratified by their regulatory outcome (e.g., increased exon inclusion \gls{vs} increased exon exclusion). Analyses can be performed for HepG2 and K562 cell lines either jointly or independently.
\\
\\
Because differential binding analyses depend critically on the \gls{dpsi} threshold used to define splicing event groups, and because the optimal threshold may vary across \glspl{rbp} due to differences in binding strength and regulatory effect size, we first conducted a systematic threshold optimisation. For each \gls{rbp} perturbation (\gls{kd}) and its corresponding binding signal derived from eCLIP data, \gls{dpsi} thresholds ranging from 0.01 to 0.1 (in increments of 0.01) were evaluated.
\\
\\
For each \gls{kd}-\gls{eclip} pair, \gls{dpsi} thresholds ranging from 0.01 to 0.10 were evaluated and the threshold maximising the statistical evidence for differential binding was selected for each splicing class and metric (see Methods, Section 9.2.5). Thresholds of 0.05 and 0.10 were additionally retained as standard reference values.
\\
\\
In the \gls{charm} interface, when users analyse a single \gls{rbp} interaction, they can select from these predefined \gls{dpsi} thresholds, which are then used to generate canonical RNA splicing maps (Figure 9.12).
\begin{figure}[H]
	\centering
	\includegraphics[width=1\linewidth]{"images/Paper Images/CHARM_Binding1"}
	\caption[CHARM's Binding Tab in Explore Mode, Single-Comparison submode]{\textbf{CHARM's binding tab in Explore Mode, single-comparison submode.} (A) Selection of splicing event type to analyse; in this example, exon skipping is selected. (B) Selection of usage mode; Explore Mode is active. (C) Selection of data scope, allowing analysis of HepG2 and K562 cell lines either jointly or separately. (D) Selection of the \gls{rbp} \gls{kd} used as the perturbation.(E) Selection of target \glspl{rbp} whose binding profiles are evaluated for changes upon the selected \gls{kd} (D). If a single \gls{rbp} is selected, results are displayed as RNA maps (see panels H–J). If multiple RBPs are selected, results are summarised as a heatmap based on the selected metric (Figure 9.11). (F) Selection of the \gls{dpsi} threshold used to define increased and decreased splicing events. The annotation in parentheses indicates the criterion by which the threshold was retained during preliminary optimisation. In this example, $dPSI = 0.10$ maximises the chi-squared statistic for increased events. (G) Selection of the metric used for visualisation. ``Effect size" corresponds to the chi-squared statistic, while ``FDR" refers to \gls{bh}\cite{benjamini_controlling_1995} adjusted p-values derived from the chi-squared test. (H) Schematic representation of the analysed splicing event type, with metagenomic coordinates provided; here, an exon skipping event is illustrated, with constitutive exons in blue, the alternative exon in red, and introns as horizontal lines. (I) Normalised binding density profile of the selected \gls{rbp} (E) under the specified \gls{kd} condition (D) and \gls{dpsi} threshold (F), computed for the splicing event category selected in (H). For each metagenome position, the normalised binding density represents the proportion of events in that category with detectable \gls{rbp} binding, excluding events with missing data at that position (see Methods, Section 9.2.4). (J) Position-wise evidence for differential binding between increased/decreased categories shown in (I) and the maintained class, represented as either the Chi-Squared statistic or \gls{fdr}-adjusted p-value according to the selection in (G). Positive values indicate positions where binding is enriched in the selected category relative to maintained events; negative values indicate depletion. The numbers of events with binding for the selected RBP is indicated between panels (I) and (J), with the total number of defined events of that type shown in parentheses. Vertical lines across (H), (I), and (J) delimitate the following metagenomic splicing regulatory sequences, in order: last 50 nt of upstream constitutive exon; first 200 nt of upstream intron; last 200 nt of upstream intron; first 50 nt of alternative exon; last 50 nt of alternative exon; first 200 nt of downstream intron; last 200 nt of downstream intron; first 50 nt of downstream constitutive exon.}
	\label{fig:Figure6.10}
\end{figure}
\noindent
When multiple interactions are explored simultaneously, the available options include the default thresholds (0.05 and 0.1) as well as four optimised thresholds, corresponding to those that maximise the selected metric (FDR or Chi-Squared statistic) for increased or decreased events within the query, represented as heatmaps (Figure 9.13).
\begin{figure}[H]
	\centering
	\includegraphics[width=1\linewidth]{"images/Paper Images/CHARM_Binding2"}
	\caption[CHARM's Binding Tab in Explore Mode, Multi-Comparison submode]{\textbf{CHARM's binding tab in Explore Mode, multi-comparison submode. }(A) Selection of usage mode; Explore Mode is active. in this example, intron retention was selected, akin to Figure 9.12A. (B) Selection of data scope, allowing analysis of HepG2 and K562 cell lines either jointly or separately. (C) Selection of the \gls{rbp} \gls{kd} used as the perturbation. (D) Selection of target \glspl{rbp} whose binding profiles are evaluated for changes upon the selected \gls{kd} (C). If a single \gls{rbp} is selected, results are displayed as RNA maps (Figure 9.10). If multiple RBPs are selected, results are summarised as a heatmap based on the selected metric (G). (E) Selection of the \gls{dpsi} threshold used to define increased and decreased splicing events. (F) Selection of the metric used for visualisation. ``Effect size" corresponds to the chi-squared statistic, while ``FDR" refers to \gls{bh}\cite{benjamini_controlling_1995} adjusted p-values derived from the chi-squared test. (G) Heatmap of position-wise evidence for differential binding between decreased (left) and increased (right) categories and the maintained category, represented as either the Chi-Squared statistic or \gls{fdr}-adjusted p-value according to the selection in (F). Positive values indicate positions where binding is enriched in the selected category relative to maintained events; negative values indicate depletion. (H) Schematic representation of the analysed splicing event type, with metagenomic coordinates provided; here, an intron retention event is illustrated, with constitutive exons in blue, and the putative retained intron in grey. Vertical lines across (G), and (H) delimitate the following metagenomic splicing regulatory sequences, in order: last 50 nt of upstream constitutive exon; first 200 nt of putative retained intron; last 200 nt of putative retained intron; first 50 nt of downstream constitutive exon.}
	\label{fig:Figure6.11}
\end{figure}
\noindent
In Discovery Mode (Figure 9.14A), users can upload an external differential splicing dataset (Figure 9.14C) to generate \gls{rna} binding maps using their own data as the source of splicing alterations, rather than the pre-loaded \gls{kd} datasets. As in the binding tab, input tables must be preprocessed using the betAS application\cite{ascensao-ferreira_betas_2024}, which standardises outputs from VAST-TOOLS\cite{tapial_atlas_2017}, ensuring compatibility with the pipeline. As in Explore Mode, users can also generate heatmaps across all available \glspl{rbp} for the event type of interest (Figure 9.14D, G). A ``Similar \glspl{rbp}" option, analogous to that described for the expression and splicing tabs, is planned for both Explore and Discovery Modes in a future release.
\begin{figure}[H]
	\centering
	\includegraphics[width=1\linewidth]{"images/Paper Images/CHARM_Binding3"}
	\caption[CHARM's Binding Tab in Discovery Mode]{\textbf{CHARM's binding tab in Discovery Mode. } (A) Selection of splicing event type to analyse; in this example, exon skipping is selected. (B) Selection of usage mode; Discovery Mode is active. (C) File upload interface, with description of file format requirements above. (D) Selection of target \glspl{rbp} whose binding profiles are evaluated for splicing changes in the user's uploaded file (C). If a single \gls{rbp} is selected, results are displayed as RNA maps (Figure 9.10). If multiple RBPs are selected, results are summarised as a heatmap based on the selected metric (G). In this example, user selected the ``All" option, which generates an heatmap for all available RBPs. (E) Selection of the \gls{dpsi} threshold used to define increased and decreased splicing events. (F) Selection of the metric used for visualisation. ``Effect size" corresponds to the chi-squared statistic, while ``FDR" refers to \gls{bh}\cite{benjamini_controlling_1995} adjusted p-values derived from the chi-squared test. (G) Heatmap of position-wise evidence for differential binding between decreased (left) and increased (right) categories and the maintained category, represented as either the Chi-Squared statistic or \gls{fdr}-adjusted p-value according to the selection in (F). Positive values indicate positions where binding is enriched in the selected category relative to maintained events; negative values indicate depletion. Vertical lines across (G) delimitate the following metagenomic splicing regulatory sequences, in order: last 50 nt of upstream constitutive exon; first 200 nt of upstream intron; last 200 nt of upstream intron; first 50 nt of alternative exon; last 50 nt of alternative exon; first 200 nt of downstream intron; last 200 nt of downstream intron; first 50 nt of downstream constitutive exon.}
	\label{fig:Figure6.12}
\end{figure}
\noindent
\\
\\
As of the writing of this thesis, the expression, splicing, and (partially) binding modules represent the currently implemented functional components of the application.
\\
\\
A planned \textbf{Network} module will integrate outputs from the preceding tabs into a unified, graph-based representation of \gls{rbp} regulatory relationships. Network edges will be derived from multiple complementary evidence layers, including shared splicing target enrichment, similarity of \gls{rna}-binding profiles, and concordance of splicing regulatory signatures across KD conditions. Together, these layers enable multi-modal inference of functional \gls{rbp}-\gls{rbp} relationships. A conceptual illustration of the planned interface is provided in Figure 9.15.

\begin{figure}[H]
	\centering
	\includegraphics[width=1\linewidth]{"images/Paper Images/CHARM_Network"}
	\caption[CHARM's Network Tab (conceptual)]{\textbf{CHARM's network tab (conceptual). } 
		(A) Selection of data scope, allowing analysis of HepG2 and K562 cell lines either jointly or separately. 
		(B) Selection of evidence layers used to construct the network. In this example, only the binding layer is active. 
		(C) When the binding layer is selected, the user can restrict the analysis to Intron Retention events, Exon Skipping events, or both. 
		(D) When the binding layer is selected, the user can further filter by directionality, retaining only increased events, only decreased events, or both. 
		(E) \Gls{mds} plot of \gls{rbp} binding-profile similarity. For each \gls{rbp} KD, the strongest eCLIP binding profile (see Section 9.2.5) was used; pairwise Euclidean distances were then computed and averaged across replicates or conditions, centring the representation on the \gls{rbp} level. As a positive control, U2AF1, U2AF2, and SF3B1 are highlighted (red ellipse): U2AF1 and U2AF2 form an obligate heterodimer, and all three proteins converge on the same intronic region at the 3$'$ splice site\cite{van_nostrand_large-scale_2020}. Their close proximity in the MDS space is consistent with their shared binding targets and supports the biological validity of the similarity metric. Whether comparable spatial proximity in this representation reflects shared regulatory function for other \gls{rbp} pairs remains to be tested experimentally.}
	\label{fig:Figure6.13}
\end{figure}